% from ega4-2-fr.tex
\oldpage[IV-2]{209}

\emph{Alg. comm.}, chap. III, § 2, n\textsuperscript{o} 13, cor. de la prop. 19).
Or, il résulte de l'hypothèse et de (7.3.15) que les fibres formelles des
$A'_{\mathfrak m_j}$ sont géométriquement normales; comme
$\operatorname{Spec}(A')$ est normal par définition, il en est de même des
$\operatorname{Spec}(A'_{\mathfrak m_j})$, et on déduit donc de (7.3.17) que
$\operatorname{Spec}(\widehat{A'_{\mathfrak m_j}})$ est normal pour tout $j$,
donc aussi $\operatorname{Spec}(\widehat{A'})$. D'autre part (Bourbaki,
\emph{Alg. comm.}, chap. IV, § 2, n\textsuperscript{o} 5, cor. 3 de la prop. 9
et chap. III, § 3, n\textsuperscript{o} 4, th. 3), $\widehat{A'}$ s'identifie
à $A'\otimes_A\widehat A$ puisque $A'$ est un $A$-module de type fini; comme
$A'$ contient $A$ et est contenu dans l'anneau total des fractions $R$ de $A$,
et que $\widehat A$ est un $A$-module plat, $\widehat{A'}$ contient
$\widehat A$ et est contenu dans $R'=R\otimes_A\widehat A$; enfin, comme
$\widehat A$ est un $A$-module plat, tout élément régulier de $A$ est aussi
$\widehat A$-régulier (0$_{\mathrm I}$, 6.3.4); donc $R'$ s'identifie
canoniquement à un sous-anneau de l'anneau total des fractions $R''$ de
$\widehat A$ (Bourbaki, \emph{Alg. comm.}, chap. II, § 2,
n\textsuperscript{o} 1, Remarque 7). Comme
$\operatorname{Spec}(\widehat{A'})$ est normal et que $\widehat{A'}$ est un
$\widehat A$-module de type fini, $\widehat{A'}$ est bien la fermeture
intégrale de $\widehat A$ dans $R''$.
\end{proof}

\begin{corollary}[7.6.2]
\label{IV.2.s7.6.2-fr}
Sous les hypothèses de (7.6.1), il y a une correspondance biunivoque canonique
entre l'ensemble des idéaux maximaux $\mathfrak m_j$ de $A'$ (autrement dit,
l'ensemble des points de $\operatorname{Spec}(A')$ au-dessus du point fermé de
$A$) et l'ensemble des idéaux premiers minimaux $\mathfrak q_j$ de
$\widehat A$ (autrement dit, l'ensemble des points maximaux de
$\operatorname{Spec}(\widehat A)$); dans cette correspondance, le complété
$\widehat{A'_{\mathfrak m_j}}$ de $A'_{\mathfrak m_j}$ est isomorphe à la
clôture intégrale de $\widehat A/\mathfrak q_j$.
\end{corollary}

\begin{proof}
On sait en effet que la fermeture intégrale de $\widehat A$ dans son anneau
total des fractions est composée directe des clôtures intégrales des
$\widehat A/\mathfrak q_j$, qui sont des anneaux locaux complets (0, 23.1.6).
\end{proof}

\begin{corollary}[7.6.3]
\label{IV.2.s7.6.3-fr}
Sous les hypothèses de (7.6.1), pour que $\widehat A$ soit intègre, il faut et
il suffit que $A$ soit unibranche; pour que $\widehat A$ soit géométriquement
unibranche, il faut et il suffit que $A$ le soit.
\end{corollary}

\begin{proof}
C'est un cas particulier de (7.6.2).
\end{proof}

\begin{theorem}[7.6.4]
\EGAReaderTarget{ega4.c04.thm.7-6-4}{IV.2.s7.6.4-fr}
(Zariski--Nagata). --- Soit $A$ un anneau semi-local noethérien. Les conditions
suivantes sont équivalentes :
\begin{enumerate}
\item[\emph{a)}] Pour toute $A$-algèbre finie réduite $C$, le complété
$\widehat C$ est un anneau réduit.
\item[\emph{a$'$)}] Pour tout anneau quotient intègre $B$ de $A$, de corps des
fractions $K$, le complété $\widehat B$ est réduit, et les corps composants
$L_i$ de l'anneau total des fractions de $\widehat B$ sont des extensions
séparables de $K$.
\item[\emph{a$''$)}] Les fibres formelles de $A$ sont géométriquement réduites
(autrement dit, pour tout anneau quotient intègre $B$ de $A$, de corps des
fractions $K$, la $K$-algèbre $\widehat B\otimes_BK$ est séparable (4.6.2)).
\item[\emph{b)}] Tout anneau quotient intègre $B$ de $A$ est un anneau
japonais.
\end{enumerate}
\end{theorem}

\begin{proof}
Pour montrer que $a)$ entraîne $a')$, il suffit de vérifier que les $L_i$ sont
des extensions séparables de $K$, ou encore (4.6.1) que pour toute extension
finie $K'$ de $K$, l'anneau $\widehat B\otimes_BK'$ est réduit; or $K'$ est
engendré par un nombre fini d'éléments entiers sur $B$, et ces derniers
engendrent une sous-$B$-algèbre finie $B'$ de $K'$, dont $K'$ est corps des
fractions. On a $\widehat{B'}=\widehat B\otimes_BB'$
((0$_{\mathrm I}$, 7.3.3) et Bourbaki, \emph{Alg. comm.}, chap. IV, § 2,
n\textsuperscript{o} 5, cor. 3 de

